\section{Finite binary descriptions of the adaptive controls}
\label{sec:digital}
The active protocol specifies a distribution by a finite command, rather
than transmitting an exact real coordinate. We separate this digital
resource from the physical production of a localized launch distribution.
The latter remains a calibrated apparatus primitive; it is not obtained
by observing the collision bits.

\begin{theorem}[A dyadic controller and its resource bounds]
\label{thm:digital}
On the same bounded class $\mathcal B_\eta$, choose the compass length
$t>0$ and the fixed coarse numerical reserves dyadic, while retaining
all strict separation margins. The experiment of
Theorem~\ref{thm:adaptive} can be implemented with dyadic nominal
centers and dyadic localization scales. Its accuracy, deterministic
attempt bound and exact discrete outputs are unchanged up to constants.
Put
\[
 b=\lceil C_0+\log_2(1/\sigma)\rceil,
 \qquad B=1+\lceil\log_2(C/\nu)\rceil
              +\lceil\log_2(C/\delta)\rceil .
\]
The finest center mesh is $2^{-b}\le c\sigma$, with
$b=O(\log(C/\nu))$. Each batch command has a binary description of
length at most $CB$, and the total transmitted command description is
at most $CJ_\nu B$. The finite boundary representation has length
at most $C\nu^{-1/(s-2)}B$. These bounds include batch repetition
counts. They do not count the analog random numbers internal to the
prescribed launch distribution as table-dependent observations.

For rational fixed numerical priors and rational $\beta$, one may
choose the fixed smooth partition and its evaluators so that a
conservative digital bit-operation bound is
\begin{equation}\label{eq:digital-work}
              C(N_\nu+J_\nu+\nu^{-1/2})B^c
\end{equation}
for an absolute finite exponent $c$. This includes boundary output,
primitive arithmetic and free-area evaluation to error $C\nu$.
It is not a sharp computational complexity bound. For general fixed
$\beta$ the message bounds hold with certified dyadic scale enclosures;
a bit-operation bound additionally depends on their computation cost.
\end{theorem}
\begin{proof}
\emph{One exact dyadic stencil.}
Take the fine localization to be a dyadic number comparable to $c h^s$,
where $h=2\pi/m\asymp\nu^{1/(s-2)}$. Choose a dyadic rectangle $W$
with all the protective collars of Lemma~\ref{lem:stopped}. Round a
requested target $x$ once to a dyadic center $x'$, with
$|x-x'|\le C2^{-b}$. Use the exact stencil
\[
                         x'+t(i,j),\qquad |i|+|j|\le n_*.
\]
For $b$ greater than the fixed denominator exponent of $t$, all stencil
points are on the same dyadic grid. Membership in $W$ is tested by
exact rational comparisons. There is no independent rounding of
transition targets, hence no error in the killed compass identity.
The reverse nominal distribution is exactly the translated joint law:
it attempts $(q+a,-a)$ for the same specified law of $(q,a)$, not an
independently rounded reverse start. Fresh attempts are still used.

The query estimates $u_\sigma(x')$. Its derived label at the intended
point $x$ is therefore correct outside the boundary tube of width
$\sigma+C2^{-b}$. Physical position, duration and angle errors receive
an independent portion of the reserve in
Proposition~\ref{prop:query}. Thus rounding and actual calibration
together preserve its fixed accuracy bound. The means are ratios of
integer hit counts to a deterministic batch size. Taking that size to
be a common integer upper bound for all queries makes every local
average and Bellman update an exact rational computation. Clipping and
comparison with $1/2$ are also exact. The depth and stencil size are
prior constants.

\emph{Rounded adaptive geometry.}
The coarse grid and its finite polygons have a prior-bounded number
of rational vertices. A dyadic interior center with a reserved fixed
inball can be found by finite search; squared distances to rational
sides test the inball inequalities. The strict inball reserve ensures
that a sufficiently fine fixed search grid succeeds. Approximate
$\cos(2\pi j/m)$ and $\sin(2\pi j/m)$ with certified error at most
$c2^{-b}$. Polygonal radial values and brackets are computed with the
same reserve. Consequently every implemented spatial midpoint is
within $C2^{-b}$ of the intended radial point.

For completeness, rounding a radial midpoint by at most $d$ does not
accumulate $kd$ after $k$ queries. The relaxed interval invariant from
Lemma~\ref{lem:bisection} holds with a boundary layer of width
$A(\sigma+d)$. Its interval widths satisfy
\[
                     w_{k+1}\le w_k/2+d,
 \qquad w_k\le2^{-k}w_0+2d.
\]
Stop before asking for resolution finer than the mesh. With
$d=C2^{-b}\le c\sigma$, $O(\log(C/\sigma))$ queries give radial
values of error $O(\sigma)$, even when labels in the enlarged layer
are chosen adaptively. Round the final values to the same mesh.
Lemma~\ref{lem:radial} absorbs these errors and still gives actual
smooth strictly convex output bodies. The known patch margin and
bounded-denominator gaps allow the same exact period decisions.

\emph{Description lengths.}
All centers lie in a fixed rectangle. Their dyadic integer numerators
therefore use $O(b)$ bits. A batch command consists of a forward/reverse
flag, two such coordinates, a scale exponent, the fixed compass/kernel
identifier, and a repetition count of order $\log(C/(\nu\delta))$.
The last count requires only $O(\log B)$ binary digits. Fixed-width
fields give the asserted $CB$ bound without a real-valued control
field. There are at most $CJ_\nu$ batches, allowing repetitions.
The receiver interprets a fixed kernel identifier as the stipulated
localized random preparation, not as a table-dependent data stream.

Store the dyadic radial values, the dyadic centers and the integers
$j,m$. In normalized coordinates the seven-point interpolation matrix
has fixed rational entries, so its coefficients have $O(B)$-bit
rational descriptions. Local angles and partitions are generated by
a fixed program from $j,m$; no table-dependent exact trigonometric
constant is stored. There are $O(m)$ coefficient blocks, proving the
boundary-description bound.

\emph{A conservative operation count.}
For the computational assertion choose an explicit smooth bump,
for example $\exp[-1/(1-x^2)]$ on $|x|<1$, zero outside, and normalize
a finite-overlap family of its translates to make the fixed partition.
The denominator has a positive fixed lower bound. To the required
precision, the bump and its first three derivatives are evaluated by
rational series on bounded arguments; near its endpoints an explicit
exponential tail bound permits certified truncation to zero. The
trigonometric functions, $\pi$, square roots and logarithms needed
here admit rational series and bisection with a polynomial number of
bit operations in $B$. For rational $s$, dyadic powers comparable to
$h^s$ can equivalently be bracketed by raising rational quantities to
fixed integer powers. Thus all scale choices are effective in this
computational subcase.

Processing bit counts, the constant-depth rational Bellman updates,
rounded bisection and the fixed-degree coefficient blocks costs
$C(N_\nu+J_\nu)B^c$. One need not optimize global numerical work to
obtain a bound: the reconstructed support functions have uniform
$C^2$ bounds. Composite trapezoidal quadrature with
$O(\nu^{-1/2})$ nodes computes their Steiner integrals, and the area
integrals $\tfrac12\int\widehat R^2$, to error $O(\nu)$. A support
value is evaluated by inverting the uniformly monotone normal-angle
map; bisection with residual enclosures uses $O(B)$ steps. Each
partition evaluation involves only a fixed number of neighboring
coefficient blocks. The uniform derivative and convexity margins
justify all these tolerances.

The number of protected components and candidate period pairs is a
prior constant. Support-defect comparisons need only a fixed precision
smaller than the known patch and rational-separation margins, whereas
the real period coordinates and areas are evaluated to $O(\nu)$.
Their mesh tests, bounded-denominator rational search and two-dimensional
Hermite arithmetic are finite with polynomial cost in $B$. Including
the preceding quadratures proves \eqref{eq:digital-work}. The larger
$\nu^{-1/2}$ term is a disclosed conservative numerical overhead; it
is not silently identified with the attempted-bit exponent.
\end{proof}

\subsection{Separate resource axes}
Write $m_\nu=\nu^{-1/(s-2)}$ and $L_\nu=\log(C/\nu)$. The following
accounting refers to the construction, not to a measured apparatus.
\begin{center}
\small
\begin{tabular}{@{}p{0.31\textwidth}p{0.64\textwidth}@{}}
\toprule
Resource & Bound or explicit status\\
\midrule
Table-dependent outputs & $N_\nu\le C m_\nu L_\nu\log(C/(\nu\delta))$ bits;
all attempts, including zeros, are charged.\\[3pt]
Command centers/batches & $J_\nu\le C m_\nu L_\nu$, counted with repetitions.\\[3pt]
Digital controls & $O(B)$ bits per batch; $O(J_\nu B)$ in total.\\[3pt]
Spatial localization & $\sigma\asymp\nu^{s/(s-2)}$; the localized kernel
is a prescribed preparation primitive.\\[3pt]
Physical calibration & $\ell+\tau+t\alpha\le c\sigma$ is sufficient.
Section~\ref{sec:calibration-floor} gives a matching necessary position-error power.\\[3pt]
Aperture & Fixed by the geometric prior, including reverse starts and
protective collars.\\[3pt]
Numerical bit operations & The conservative bound \eqref{eq:digital-work}
under its explicit effective-prior assumptions.\\[3pt]
Motion, setup, calibration production & No physical time, energy or
cost law is assumed. No bound on these resources is inferred from $N_\nu$.\\
\bottomrule
\end{tabular}
\end{center}
Digital description length is not mechanical resolution. In particular,
a finite word specifying a narrow kernel does not prove that a device
can produce that kernel cheaply. The converse calibration theorem
below quantifies a necessary tolerance in the observation model,
without inventing a law for the cost of achieving it.
